\section{Implications for how cough detectors are evaluated}\label{sec:discussion}

The results above are specific to three corpora and one pipeline, but the way they fail
generalises, and the practices that would have exposed the failure earlier are neither
expensive nor novel. This section sets out what follows for anyone building or reviewing
a detector of this kind. Section~\ref{sec:disc-protocol} concerns the evaluation protocol,
Section~\ref{sec:disc-regime} concerns what a reported figure means once its label regime
is known, and Section~\ref{sec:disc-corpora} concerns what the corpora themselves would
need to contain.

\subsection{What an adequate protocol costs}\label{sec:disc-protocol}

The external evaluation reported here consumed a few minutes of inference over $476$
recordings, on weights that already existed, against a corpus nobody trained on. It
required no annotation, no collection and no additional training. Set against the cost of
the grid that produced those weights, it is a rounding error, and it is the only step in
this study that changed what could be concluded. The audits of this
literature~\cite{roberts2021, han2022} record that external validation is usually absent;
the present results suggest that its absence is not explained by its cost.

Four reporting practices follow directly from what the numbers here show, and each
addresses something a single accuracy figure concealed.

The first is to report sensitivity and specificity separately rather than, or at least
alongside, any quantity that averages them. The fall described in
Section~\ref{sec:res-external} is almost entirely a fall in sensitivity, and an
application that monitors a patient for cough frequency is harmed by exactly that
component. Balanced accuracy halves the visible damage, and plain accuracy on a prior of
this shape hides it almost completely.

The second is to report an estimate of run to run variation. Every condition here was run
under five seeds, and the spread that produced is the quantity against which an effect has
to be judged. The denoising effects in Section~\ref{sec:res-denoise} are smaller than that
spread in three of the comparisons and comparable to it in the fourth, so none of them
could have been claimed from a single run however large the test partition was. A
difference of a few tenths of a point, reported without a spread, is not a result.

The third is to report the area under the curve beside the thresholded figures. It cost
nothing here, since the scores were already being saved, and it answered a question the
thresholded numbers could not: whether a practitioner facing this degradation should
reach for recalibration. Table~\ref{tab:threshold} says they should not, and without the
ranking statistic there would have been no way to tell.

The fourth is to report what the false positives are, and not merely how many. The
aggregate specificity of $84.3$ per cent in Table~\ref{tab:external} is not an alarming
number on its own. Table~\ref{tab:confusion} shows that it is composed of a severe and
localised failure on sneezes together with much lower rates elsewhere, and it is the
localisation, not the aggregate, that identifies what is missing from the training data.

A fifth practice concerns the partition rather than the reporting. The partition here is
built over participants, which is what makes it possible to say that the in corpus figures
are not speaker recognition in disguise. That was not sufficient to make them transfer,
which is the point of this paper, but it was necessary in order to interpret them at all.
A corpus that does not publish a participant identifier cannot support this, and
Section~\ref{sec:res-regime} describes one that does not.

\subsection{What a figure means once its regime is known}\label{sec:disc-regime}

The three label regimes were evaluated with one pipeline, one architecture and one input
condition held fixed across them, which makes the three resulting figures comparable in a
way that figures drawn from three papers are not. A ResNet18 on linear spectrograms
reaches $98.60 \pm 0.22$ per cent balanced accuracy over held out speakers whose labels
came from what they were asked to do, $98.41 \pm 0.24$ where the labels came from another
classifier's score, and $70.18 \pm 0.44$ where they came from a human listening to the
recording. All three would be reported as cough detection accuracy. The first two are
indistinguishable at the precision anyone reports, and the third sits almost thirty points
below either.

Set against that, the effect of every modelling choice measured here is small. On linear
spectrograms the widest gap between architectures separates the four layer network's
$95.33 \pm 0.57$ from the same ResNet18's $98.60$, which is about three points; the widest
effect of the frequency axis, holding the architecture fixed, is under half a point; and
wavelet shrinkage, the intervention the original work credited with its improvement,
detectably improves none of the nine paired comparisons made across the two trainable
corpora, the four in which its effect is detectable at all being losses. The provenance of the label
therefore dominates the design of the model by roughly an order of magnitude. A
practitioner choosing where to spend effort
should read that ordering literally: a better architecture applied to prompt derived
labels will not recover what the labels themselves give away.

This has a consequence for reading the literature that is worth stating bluntly. Reported
accuracies for this task are not on a common scale, and the quantity that moves them most
is the one least often named. Two figures are only comparable when both corpora obtained
their labels the same way, both splits were built on the same kind of identifier, and both
priors match; absent any of those, the difference between two published numbers carries no
information about the two detectors. The remedy is cheap and entirely within an author's
control. Name the regime beside the number, say which identifier the split was built on,
and report one corpus whose labels were produced by a different process from the training
corpus. The third of those is the only one that costs anything, and
Section~\ref{sec:disc-protocol} has already put a figure on what it costs.

\subsection{What the corpora would need to contain}\label{sec:disc-corpora}

The concentration of errors described in Section~\ref{sec:res-external} points at a
property of the training corpus rather than of the model. A detector trained against
breathing, sustained vowels and counting is never required to separate coughing from any
other forced expiration, and it turns out not to. The remedy is not more data of the same
kind. A corpus ten times the size, collected under the same protocol, would supply ten
times as many easy negatives and would leave the distinction that matters untested.

What would help is a negative class assembled against the confusions the task actually
presents. Sneezing, throat clearing, burping, gasping and sighing all share the mechanism
that makes a cough acoustically distinctive, and a prompted collection could obtain them
the same way it obtains coughs, by asking. That the human annotated corpus used here
supplies such material at all is a consequence of its ontology rather than of any
intention to support this task, and the $106$ cough clips it contains are too few to
train on.

Label provenance should also be stated rather than assumed. The three regimes examined
here are routinely treated as interchangeable sources of ground truth, and they support
different claims. A prompt establishes what a participant was asked to do. An annotation
establishes what a listener judged to be present. A detector score establishes what
another model believed, and a figure obtained by training and testing against it measures
agreement between two classifiers, which is not a clinical quantity at all. None of the
three is useless, and the first in particular is what makes a matched speaker disjoint
partition possible here. The difficulty is in reporting a number obtained under one of
them as though it had been obtained under another.

Finally, a corpus that publishes a derived label should publish the procedure that
produced it in a form that can be re-executed. Section~\ref{sec:res-regime} documents a
selection that cannot be reproduced from the corpus as released, because the reported
count falls strictly between two consecutive attainable values, and an expert annotated
subset that cannot be used to estimate the corpus error rate because it was drawn from
where the detector was already confident. Neither property is a defect of the recordings,
and both limit what any study built on them can establish.
